\documentclass[11pt]{article}
\usepackage[margin=0.9in]{geometry}
\usepackage{amsmath,amssymb,amsthm,booktabs,mathtools}
\usepackage[hidelinks]{hyperref}
\usepackage{pgfplots,needspace}
\pgfplotsset{compat=1.18}
\newtheorem{theorem}{Theorem}
\newtheorem{proposition}{Proposition}
\newcommand{\T}{\mathbb T}
\newcommand{\N}{\mathbb N}
\newcommand{\Z}{\mathbb Z}
\newcommand{\HS}[1]{H^{#1}(\T)}
\title{Recovery of Polarized Gowdy Fields from Two Time Slices}
\author{Ronald Bibb}
\date{Working draft, 25 November 2026}
\begin{document}
\maketitle

\begin{abstract}
We consider two full spatial snapshots of a polarized Gowdy field on
a marked torus at finite areal times. For a rational separation
$b-a=\pi p/q$ and outside a locally finite set of base times, the
two-snapshot map determines every smooth field solution. Its stability
depends on the target. Field slices on the lattice $a+(\pi/q)\Z$
can be recovered without loss of Sobolev regularity. Every other fixed
positive target time costs exactly one derivative. Recovery of the limiting
log-time slope at the singular end costs exactly three halves.
The distinction comes from an order-$k^{-1}$ shift of the high-frequency
phase. This shift lifts the leading-order resonances, while the target
slice can cancel their contribution. Numerical checks compare the exact
Bessel transfer with an independent ODE integration. Sampling each
frequency band at its own time gives uniform conditioning with two
stored coefficients per mode and an order-logarithmic number of
observation times. A finite-window test in nonlinear vacuum Gowdy
illustrates the effect of the reconstruction rule. For the original
production case, the localized score falls from $0.822$ to $0.00298$
using the same observations and a different restart.
\end{abstract}

\section{Two snapshots and the choice of target}
The polarized Gowdy equations in areal time $t>0$ include
\begin{equation}
 f_{tt}+t^{-1}f_t-f_{\theta\theta}=0,\qquad \theta\in\T=\mathbb R/(2\pi\Z).
 \label{wave}
\end{equation}
Here $f$ is the polarized wave field $P$, or an identically evolving
rescaled massless scalar in the symmetric Einstein--scalar model. We hold
the coordinate period and the marking fixed. Two measurements of $f$
can be encoded by a secant slope and an intercept. The conditioning of
recovery depends on the evolution represented by those measurements.
A large error from a prescribed restart rule alone leaves both the
injectivity and the stability of the underlying observation map open.

For constant-coefficient waves, pointwise-in-time observation theory
relates stability to arithmetic properties of the sampling times. Vest's
loaded-string result treats two rationally spaced measurements with
a one-derivative loss \cite{szi,vest}. Christensen et al. identify
a lattice of determined snapshots for the free wave equation on
$\mathbb R^n$ \cite[Theorems 2.6--2.7]{snapshot}. For Gowdy,
Amorim, Bernardi and LeFloch give two numerical evolution schemes,
examine sharp profiles, and propose a Cauchy-to-asymptotic-to-Cauchy
strategy \cite[Sections 3 and 4.2.4--4.2.5]{amorim}. We ask which
specific targets are stable under the same finite-time observation map.
The result uses a phase correction from the $t^{-1}f_t$ term and a
comparison of its effect at the observation and target times.
It gives sharp estimates for the field and the limiting slope on the same
data. Li's Kasner scattering estimates describe the regularity of singular
data when full Cauchy data are given \cite[Theorem 1.3]{li}. In our
two-snapshot problem, recovering a missing Cauchy component contributes
an additional loss.

\section{Sharp finite- and infinite-depth estimates}
Set $y=\sqrt t\,f$ and $W(t)=1/(4t^2)$. For each nonzero Fourier mode
write $k$ for the magnitude of the mode number. Then
\begin{equation}
 y_k''+(k^2+W)y_k=0 .
 \label{ode}
\end{equation}
Let $C_k(t)$ and $S_k(t)$ have data $(1,0)$ and $(0,k)$, respectively,
at $a$. The exact snapshot formulas are
\begin{equation}
 \frac{y_k'(a)}{k}=\frac{y_k(b)-C_k(b)y_k(a)}{S_k(b)},\qquad
 y_k(c)=C_k(c)y_k(a)+\frac{S_k(c)}{S_k(b)}
 [y_k(b)-C_k(b)y_k(a)].
 \label{multipliers}
\end{equation}
The zero mode of $f$ is $A+B\log t$ and is determined at any two
distinct positive times. Let $M_s=\|f(a)\|_{\HS{s}}+\|f(b)\|_{\HS{s}}$.
The estimates below also apply to differences of solutions.

\begin{theorem}
\label{thm}
Fix $L=\pi p/q>0$ with $p,q$ positive and coprime. For $a>0$ set
$b=a+L$. There is a locally finite set of exceptional $a$.
Outside this set, the two snapshots determine a unique smooth solution.
A sufficient condition for exclusion from the exceptional set is
$\int_a^b W(u)\,du<\pi/q$. The alternative sufficient bound $\int_a^b W(u)\,du<\pi/(p+1)$ permits pairs arbitrarily near zero by reducing the gap.
For every real $s$ the following hold.
\begin{enumerate}
\item The Cauchy field velocity $f_t(a)$ has
$\|f_t(a)\|_{\HS{s-2}}\lesssim M_s$ with sharp loss two for the unscaled
velocity. Equivalently, relative to the usual order-one differentiation
in the Cauchy-data norm, reconstruction costs one additional derivative.
\item For fixed $c>0$, if $q(c-a)/\pi$ is an integer, then
$\|f(c)\|_{\HS{s}}\lesssim M_s$. Otherwise
$\|f(c)\|_{\HS{s-1}}\lesssim M_s$, and every smaller loss fails.
\item For smooth solutions let $v=-\lim_{t\downarrow0}t f_t$.
Then $\|v\|_{\HS{s-3/2}}\lesssim M_s$, and every smaller loss fails.
\end{enumerate}
The constants can depend on the fixed times. The estimates with
Sobolev data extend from smooth data by continuity.
\end{theorem}

\begin{proof}
On each compact time interval in $(0,\infty)$ containing $a$, a
Volterra representation gives
\begin{align}
 C_k(t)&=\cos(k(t-a))+O(k^{-1}),\\
 S_k(t)&=\sin(k(t-a))
       +\frac{\cos(k(t-a))}{2k}\int_a^t W(u)\,du+O(k^{-2}).
 \label{phase}
\end{align}
For the second equation, substitute the leading sine once into
$S_k(t)=\sin(k(t-a))-k^{-1}\int_a^t \sin(k(t-u))W(u)S_k(u)\,du$.
The product-to-sum formula gives the displayed correction. Integration
by parts controls the oscillatory term. A Volterra bound controls the
second and higher iterations. These estimates are uniform for $a,b$ in
a compact subset of the positive times. For $q\nmid k$, the leading
sine at $b$ is bounded away from zero for $q>1$. For $q\mid k$, the
first correction is nonzero and $|S_k(b)|\asymp k^{-1}$.

For the sufficient condition, write
$S_k=r_k\sin\varphi_k$,
$S_k'/k=r_k\cos\varphi_k$, with $r_k(a)=1$ and $\varphi_k(a)=0$.
Then $\varphi_k'=k+(W/k)\sin^2\varphi_k$ and $r_k>0$.
The phase at $b$ is strictly greater than $kL$ by at most
$k^{-1}\int_a^b W(u)\,du$. The sufficient bound places that phase
strictly below the next multiple of $\pi$, including when $kL$
starts at a multiple. For the alternative bound set $I=\int_a^b W$.
The first argument applies if $I/k<\pi/q$. In the remaining modes
$kL+I/k\le pI+I<\pi$, with positive phase. For general $a$, the large-$k$ bounds exclude
all but finitely many modes on each compact subinterval of $a>0$.
For each remaining mode $S_k(a+L)$ is analytic as a function of $a$
by the Bessel formula for \eqref{ode}. It is nonzero for sufficiently
large $a$ by the sufficient condition. Its zeros are isolated. These
facts give local finiteness of the exceptional set. The inverse
of \eqref{multipliers} has polynomial growth outside this set.
This also yields the smooth-data existence statement for arbitrary
smooth snapshot pairs.

The first part follows from \eqref{multipliers}. For sharpness take
one mode at a time with $y_k(a)=0$, $y_k'(a)=k$ and $k$ divisible by
$q$. Its sole nonzero snapshot amplitude is of order $k^{-1}$,
and its Cauchy velocity amplitude is of order $k$.
These real single-mode fields satisfy the polarized momentum
integrability condition when taken as vacuum $P$ data.
For the second part use $S_k(c)/S_k(b)$ in \eqref{multipliers}.
It is bounded for $q\nmid k$. For $q\mid k$ the quotient is
bounded when $q(c-a)/\pi\in\Z$, since both leading sines vanish.
At any other $c$ there are infinitely many $k=qm$ with
$|\sin(k(c-a))|$ bounded below. For rational noninteger
$q(c-a)/\pi$ take a nonzero residue class. For irrational values use
density of the irrational rotation. Along that sequence the same
single-mode examples have target amplitude bounded below and
snapshot amplitude $O(k^{-1})$. Multiplication by the constant
factors $\sqrt a,\sqrt b,\sqrt c$ leaves all exponents unchanged.

For the last part write
$f_k(t)=A_kJ_0(kt)+B_kY_0(kt)$.
The Wronskian is $2/(\pi t)$ in areal time. Thus
\[
 B_k=\frac{\pi}{2}\sqrt a\left[
 J_0(ka)y'_k(a)
 -\left(kJ'_0(ka)+\frac{J_0(ka)}{2a}\right)y_k(a)\right].
\]
The slope is $-2B_k/\pi$. The small-argument limits and high-frequency
amplitudes of $J_0$ and its derivative imply
$|v_k|\lesssim \sqrt k (|y_k(a)|+|y'_k(a)|/k)$
\cite[Sections 10.5, 10.7 and 10.17]{dlmf}. This gives the upper
bound. On the single-mode tests $|v_k|=k\sqrt a |J_0(ka)|$.
Along an infinite subsequence of $k=qm$, $ka$ approaches $0$ modulo
$\pi$, or equals $0$ modulo $\pi$ when $qa/\pi$ is rational.
The leading phase of $J_0$ is $ka-\pi/4$, so the slope is bounded
below by a constant times $\sqrt k$ along that subsequence.
The ratio to the snapshot amplitude is at least order $k^{3/2}$.
The zero mode obeys finite-dimensional bounds. Modewise convergence
defining $v$ for smooth data is consistent with the smooth polarized
asymptotics \cite{cim}. This completes the proof.
\end{proof}

The finite-time target result also holds for the transformed periodic
wave equation with a time-only $C^1$ potential on the compact interval
of interest, provided its signed integral between the observations
is nonzero and every low-mode denominator in \eqref{multipliers}
is nonzero. Indeed these are the hypotheses used for the quotient
bounds and the off-lattice counterexamples. A strictly positive
potential and the small-integral condition give a uniform way to
verify the low-mode hypothesis. The limiting slope claim uses the
inverse-square Gowdy potential.

Actual collisions can occur when the observation multiplier vanishes. For fixed $a$ and a nonzero mode $k$, the
nontrivial $S_k$ has positive-time zeros beyond $a$ by Sturm comparison
with the constant coefficient $k^2$. Select one such $b$.
The polarized vacuum fields $P=0$ and
$P=\epsilon S_k(t)\cos(k\theta)/\sqrt t$ agree at $a,b$
and differ at intervening slices. Both satisfy the periodic
momentum constraint. Each solution's free metric constant
can be chosen to give the same single metric anchor at $a$.
One history then has a spatially varying $P$ at the target,
and the other has constant $P$. This supplies a case of actual
non-injectivity for selected finite snapshots and a different sampling
regime from Theorem~\ref{thm}.

\paragraph{Marked geometry at a stable field slice.}
The field $P$ gives half the logarithm of the ratio of the squared
lengths of the two marked Killing generators on a polarized slice.
Stability of this field target need not extend to its normal
derivative. In vacuum the directional expansion asymmetry is
$s_P=(H_x-H_y)/\Theta=-4P_\tau/(3+\lambda_\tau)$ in the convention
below. Linearize about $P=0$, $\lambda=0$ with fixed metric anchor.
The first variation is $-4\delta P_\tau/3$.
For a lattice time $c$, the same single-mode input used above has
$|S_k'(c)|\asymp k$ on $q\mid k$, by the differentiated Volterra
formula. Thus the linearized $s_P$ costs exactly two derivatives from
two field snapshots even when the intrinsic field $P(c)$ has no loss.
The upper estimate follows from $|S_k'|+|C_k'|=O(k)$ and the inverse
multiplier. This claim concerns the derivative of the observable map
at a homogeneous background.

\Needspace{6\baselineskip}
\begin{proposition}[Timing tolerance at a finite bandwidth]
Under the nonvanishing hypothesis of Theorem~\ref{thm}, consider a
spectral cutoff $K$ and an on-lattice target $c>0$. Allow the known
actual observation time and target to be $b+\epsilon$ and $c+\eta$.
There are fixed $\epsilon_0,\eta_0>0$ such that the target gain remains
bounded independently of $K$ for $|\epsilon|\le\epsilon_0 K^{-2}$
and $|\eta|\le\eta_0 K^{-2}$.
The scale $K^{-2}$ is sharp in two senses.
\begin{enumerate}
\item For all sufficiently large $k$ divisible by $q$ there is a blind
observation time
\[
 b_k=b-\frac{1}{2k^2}\int_a^b W(u)\,du+O(k^{-3}).
\]
Allowing two-sided observation perturbations on a scale asymptotically
larger than $K^{-2}$ eventually includes a kernel at some $k\le K$.
\item With observation times fixed, a target-time tolerance $r_K$
such that $K^2r_K\to\infty$ gives unbounded worst-case target gain on
$|\eta|\le r_K$.
\end{enumerate}
\end{proposition}
\begin{proof}
On fixed compact intervals $|C_k|\lesssim1$ and
$|S_k'|\lesssim k$ by differentiation of the Volterra formula.
For $k\le K$ the change in $S_k(b)$ is at most
$O(k|\epsilon|)\le O(\epsilon_0/k)$, and
$|S_k(b)|\gtrsim1/k$ on all modes with $k\ge1$.
Choose $\epsilon_0$ to retain half of this lower bound.
On multiples of $q$, $S_k(c)=O(k^{-1})$, and this remains true
when $|\eta|\le\eta_0K^{-2}$. On nonmultiples the denominator
starts uniformly bounded below and can also retain half of that bound.
Use the multiplier formula and the explicit zero mode to obtain
the uniform estimate. All perturbations are small enough to keep
times positive and the observation pair distinct.

For the first sharpness statement, $S_k(b)=(-1)^{kp/q}
(\int_a^b W/(2k))+O(k^{-2})$ when $q\mid k$.
For $|t-b|=O(k^{-2})$ the differentiated Volterra formula yields
$S_k'(t)=(-1)^{kp/q} k+O(1)$. The intermediate value theorem
and this derivative give the unique nearby root and its expansion.
Taking the largest multiple of $q$ at most $K$ gives the stated
obstruction. For the second sharpness statement, use that same mode
and take $\eta=\min(r_K,1/(2k))$, further capped if needed to stay
on a fixed compact positive interval. The leading sine at $c+\eta$
has absolute value $\gtrsim k\eta$. Dividing by $S_k(b)$ gives
growth at least a constant times $\min(k^2 r_K,k)$ up to a bounded
error. The multiplier of the second snapshot alone is that quotient.
\end{proof}

\section{A frequency-dependent storage rule}
\label{schedule}
The arithmetic obstruction comes from storing every mode at the
same second time. For a simulation with access to Fourier coefficients
during evolution, keep one coefficient at the reference slice and
one at a band-dependent time. Each nonzero mode uses the same number
of stored coefficients as in two complete snapshots.

\begin{proposition}[Band-dependent observation times]
Fix $a>0$, a compact right observation interval $[a,a+T]$, and a
frequency cutoff $K$. There is a schedule using
$1+\lfloor\log_2 K\rfloor$ distinct additional times with uniform
reconstruction of $(y_k(a),y_k'(a)/k)$ for $1\le k\le K$.
Each mode uses one additional measurement.
The order $\log K$ is necessary for uniform modewise conditioning
with the maximum norm on the measured coefficients, even if all modes
are available at each time and the times may depend on $K$.
\end{proposition}

\begin{proof}
Let $W_* = \sup_{[a,a+T]} W$. Choose $0<d\le
\min\{T,\pi/[4(2+W_*)]\}$. For $2^j\le k<2^{j+1}$
save the second coefficient at $t_j=a+d\,2^{-j}$.
The Pr\"ufer phase there lies between $d$ and $(2+W_*)d\le\pi/4$.
The amplitude is bounded below uniformly because
$|r_k'/r_k|\le W_*/(2k)$.
Consequently $S_k(t_j)$ is bounded away from zero uniformly.
The zero mode uses $t_0$ and its elementary two-time formula.
Solving \eqref{multipliers} with $b=t_j$ gives the conclusion.
With the natural Sobolev norm on the sequence of stored coefficients,
every fixed positive-time field is recovered without loss, and the
limiting slope still costs one half derivative from its high-frequency
Bessel factor. The single-mode tests with $ka$ near $0$ modulo $\pi$
show sharpness of the half derivative.

For the lower bound, suppose that an estimate with constant $M$,
independent of $K$, holds for the missing normalized velocity in the
maximum norm of the observed coefficients. Setting the reference
coefficient to zero and normalized velocity to one gives
$1\le M\max_i |S_k(t_i)|$ for all $1\le k\le K$.
On the fixed time interval, \eqref{phase} gives an error bounded by
$B/k$ in replacing $S_k(t_i)$ by $\sin(k(t_i-a))$.
Fix an integer $\ell> 4MB$ with $\ell\ge1$ and an integer
$R>4\pi M$ with $R\ge2$.
Partition the unit $m$-cube into $R^m$ equal boxes and compare
$R^m+1$ multiples of the vector with entries
$\ell(t_i-a)/\pi$ modulo one. Two inhabit the same box.
Their positive index difference gives a mode
$\ell\le k\le \ell R^m$ with $|\sin(k(t_i-a))|\le\pi/R$
for every $i$. If $\ell R^m\le K$ this contradicts the uniform
bound. Thus $m>\log_R(K/\ell)$.
\end{proof}

The schedule assumes full spatial data at the chosen times or
the ability to compute a Fourier band in situ. It stores as many
coefficients as two complete snapshots, while requiring access at
more distinct times. Perturbing a band time by at most a sufficiently
small multiple of $1/K$ retains the uniform bound through cutoff $K$,
since $S_k'=O(k)$. This tradeoff is different from observability
using a fixed number of time instants and a weakened Cauchy norm
\cite[Theorem 4.1]{vest}. It is related to dynamical sampling, which
combines measurements at different times under a specified evolution
\cite{aldroubi}.

\paragraph{Euler--Poisson--Darboux extension.}
The finite-time result and the storage rule extend to
$f_{tt}+(\alpha/t)f_t-f_{\theta\theta}=0$ for $0<\alpha<2$.
Indeed $y=t^{\alpha/2}f$ gives
$W=\alpha(2-\alpha)/(4t^2)>0$. The same analysis gives a sharp
snapshot-to-singular-flux loss $2-\alpha/2$ with a rational pair,
and $1-\alpha/2$ with the band schedule. Here the target flux is
$\lim_{t\downarrow0}t^\alpha f_t$.
To see the new exponent, put $\nu=(\alpha-1)/2$ and use
$f_k=t^{-\nu}(A_kJ_\nu(kt)+B_kY_\nu(kt))$.
For $\nu\ne0$ the term proportional to $t^{-2\nu}$ in $f_k$
has coefficient a nonzero constant times $k^{-\nu}B_k$, by the
$J_{-\nu}$ term in the connection formula for $Y_\nu$. The
constant term in $f_k$ has zero limiting flux. The $\nu=0$
calculation gives the logarithmic case above.
The Wronskian for $\sqrt tJ_\nu(kt),\sqrt tY_\nu(kt)$ gives
$|B_k|\lesssim k^{1/2}(|y_k(a)|+|y_k'(a)|/k)$. On the
single-mode witnesses with $y_k(a)=0$ and $y_k'(a)=k$, the
leading Bessel phase gives $|B_k|\gtrsim\sqrt k$ along the
subsequence with $ka\to0$ modulo $\pi$. The temporal inversion
factors then give the claimed losses. At $k=0$ integrate the
power-law solution directly. Sharp regularity estimates for
the singular Cauchy problem for this PDE predate the present
observation problem \cite{epd}.

\section{Numerical transfer check}
Choose $a=\pi/4$ and $b=a+3\pi/16$. The target $c=a+4\pi/16$
lies on the stable lattice; $c'=a+3.5\pi/16$ lies off it.
For observations of $y$ the Euclidean gain to one $y$ coefficient is
\[
 G_k(c)=\left(\left|C_k(c)-C_k(b)S_k(c)/S_k(b)\right|^2+
                      \left|S_k(c)/S_k(b)\right|^2\right)^{1/2}.
\]
On modes $1$ through $65536$ the largest sampled gain at $c$ was
$6.04$. At $c'$ it reached $1.36\times10^6$ at mode $65520$,
on the predicted linear-growth subsequence. On near-resonant modes
$|kS_k(b)|$ approaches $\frac12\int_a^b W=0.0682093$.
Six comparisons between the exact Bessel formula and a separate
DOP853 integration gave maximum error $2.51\times10^{-12}$ in
the two transfer components. These computations illustrate
the continuum argument, whose bounds and lower witnesses follow
from \eqref{phase}. Shifting the second observation to a modal
root gave $k^2(b-b_k)=0.0682023$ at $k=128$ and $0.06820925$
at $k=8192$. Moving the lattice target by $1/k^2$ gave gain
close to $21.69$ on those tested modes. Moving it by $1/(2k)$
raised the gains from about $1.27\times10^3$ to $8.14\times10^4$.
The band schedule used 17 additional times, with one coefficient
from those times per mode. Its maximum sampled normalized
velocity gain was 4.29. Figure~\ref{fig:gains} compares the
field gain at the two target slices.
\begin{figure}[ht]
\centering
\begin{tikzpicture}
\begin{loglogaxis}[
  width=0.92\linewidth,height=6.2cm,
  xlabel={Spatial Fourier mode $k$},ylabel={Gain to the weighted field $y$},
  grid=major,major grid style={gray!20},
  legend cell align=left,legend pos=north west,
  tick label style={font=\small},legend style={font=\footnotesize,draw=none}
]
\addplot+[red!70!black,mark=triangle*,thick] table[x=k,y=fixed_offset,col sep=comma]{mode_gains.csv};
\addlegendentry{Fixed pair, off-lattice target}
\addplot+[blue!70!black,mark=*,thick] table[x=k,y=fixed_aligned,col sep=comma]{mode_gains.csv};
\addlegendentry{Fixed pair, lattice target}
\addplot+[teal!75!black,mark=square*,dashed,thick] table[x=k,y=band_offset,col sep=comma]{mode_gains.csv};
\addlegendentry{Band schedule, off-lattice target}
\end{loglogaxis}
\end{tikzpicture}
\caption{Selected modes $k=16(2m+1)$ on the near-resonant subsequence.
The two targets differ in areal time by $\pi/32$.
The band schedule uses the same number of stored coefficients per mode.}
\label{fig:gains}
\end{figure}

\paragraph{A late-window check.}
For comparison with the original extraction window, use backward log
time $u=\tau_f-\tau$ and let $L_k(u)$ be the unit initial derivative
solution of $L_k''+k^2e^{-2(\tau_f-u)}L_k=0$ with $L_k(0)=0$.
If $M=K^2e^{-2(\tau_f-\Delta)}$ and $\Delta\sqrt M<\pi$,
Sturm comparison gives $L_k>0$ on $(0,\Delta]$ for $k\le K$.
Concavity gives $L_k(u)\le u$. The Volterra formula then gives
\[
  1-M\Delta^2/6\ \le\ L_k(\Delta)/\Delta\ \le 1.
\]
For $\Delta=2$ and the Nyquist bound $K=2048$, the relative bound
$M\Delta^2/6$ is $5.8\times10^{-3}$ at $\tau_f=12$ and
$3.6\times10^{-8}$ at $\tau_f=18$. These are statements about the
polarized continuum modal evolution over the recorded window.
They exclude a linear phase ambiguity through this bandwidth
for those windows. The nonlinear coupling can still give a
substantial bias to a slope-based velocity estimate.

\section{Corrected polarized Einstein--scalar benchmark}
\label{metric}
Write $\tau=-\log t$, take $8\pi G=1$, and let $\varphi$ denote
the canonical massless scalar. Use $\phi=\sqrt{2}\varphi$ for the
equal-coefficient field variable. With positive metric quadratures
the metric is
\begin{align}
 ds^2={}&e^{(\tau-\lambda)/2}(-e^{-2\tau}d\tau^2+d\theta^2)\notag\\
   &+e^{-\tau}(e^Pdx^2+e^{-P}dy^2),\\
 \lambda_\tau={}& P_\tau^2+\phi_\tau^2
                 +e^{-2\tau}(P_\theta^2+\phi_\theta^2),\\
 \lambda_\theta={}& 2(P_\tau P_\theta+\phi_\tau\phi_\theta).
 \label{constraint}
\end{align}
The original preprint used the opposite sign for $\lambda$ in the
metric and a different scalar density normalization \cite{earlier}.
In the convention with $e^{(\tau+\lambda)/2}$, as in Lim,
both quadratures change sign \cite[Eqs.~(1), (4)--(5)]{lim}.
With normal pointing toward increasing $\tau$, the lapse and
geometric rates are
\begin{align}
 N&=e^{(-\lambda-3\tau)/4},\qquad
 H_\theta=\frac{1-\lambda_\tau}{4N},\\
 H_x&=\frac{P_\tau-1}{2N},\qquad H_y=\frac{-P_\tau-1}{2N},\\
 \Theta&=-\frac{3+\lambda_\tau}{4N},\\
 \rho&=\frac{\phi_\tau^2+e^{-2\tau}\phi_\theta^2}{4N^2}.
\end{align}
Let $H=\Theta/3$, $\Omega=\rho/(3H^2)$,
$\Sigma^2=\sum_i(H_i-H)^2/(6H^2)$, and $s_P=(H_x-H_y)/\Theta$.
The comparison uses these quantities on equal areal slices and the
dimensionless rate $\Theta e^{-\tau/2}$. A separate relative difference
in $\rho^{-1/2}$ compares density scales on the same slices.
If the unit normal is reversed, $\Theta$ changes sign and the
three normalized shape quantities stay fixed.

Among allowed asymptotic fields $P=v_P\tau+P_\infty+o(1)$,
$\phi=v_\phi\tau+\phi_\infty+o(1)$, assume differentiable asymptotics
of sufficient order to integrate the quadratures. The periodic
compatibility condition is
$\int_{\T}(v_PP_\infty'+v_\phi\phi_\infty')\,d\theta=0$.
With $q_0=v_P^2+v_\phi^2$, one has
$\lambda=q_0\tau+\lambda_\infty+o(1)$ and
$\lambda_\infty'=2(v_PP_\infty'+v_\phi\phi_\infty')$.
The integration constant $c_\lambda=\lambda_\infty(0)$ is
independent on the fixed marked torus \cite[Eq.~(3.19)]{beyer}. The shift
$\lambda\mapsto\lambda+c$ rescales $N$ and the transverse proper
length by $e^{-c/4}$, and normal rates by $e^{c/4}$.
These factors cancel from the normalized shape score.
The score also misses a constant shift of $P$, which changes
the two marked orbit lengths in reciprocal directions.
The full-state reconstruction of the marked metric requires
retaining its actual field data and integration constant.

The smooth production example has
$P=0.6\cos\theta$, $P_\tau=0.4+0.3\cos 2\theta$,
$\phi=0.5\sin\theta$, and $\phi_\tau=0.8-0.2\cos 2\theta$
at $\tau=0$, with $\lambda(0,0)=0$.
The reconstruction fits the two fields from the slices
$\tau_f-4$ and $\tau_f-2$ and estimates
$c_\lambda$ from $\lambda(\tau_f-2,0)-\widehat q_0(0)(\tau_f-2)$.
This supplies one pre-endpoint metric measurement in addition to
the field samples. Truth at the restart surface is withheld
from the main estimator. Centered second-order spatial differences
and RK4 evolve both branches.

At $N=1024$ the maximum geometric comparison distances on
$\tau=2,4$, for $\tau_f=8,12,16,24$, were
$4.759\times10^{-4}$, $1.060\times10^{-5}$,
$9.786\times10^{-6}$, and $9.786\times10^{-6}$ respectively.
At depths 16 and 24 each resolution doubling from 256 to 512 and
1024 lowered the error by about four. At depth 8 the refinement
sequence approached a finite-window error floor. Setting the estimated
constant to zero while retaining the four field functions produces
a score of about $0.232$ at the finest grid across all four depths.
In the analytic multimode Bessel benchmark the depth-16 fine-grid
scores were $1.3344\times10^{-5}$ both when the limiting coefficients
were supplied and when pre-endpoint holdout was used. The analytic
arm's finite-time metric truth uses the constrained quadrature of
the known fields. These results verify the declared outputs in the
chosen family and areal slicing.

\section{A finite-window nonlinear inverse}
\label{nonlinear}
For vacuum unpolarized Gowdy let $V=P_\tau$ and $R=Q_\tau$.
Given $P(f-\Delta),P(f),Q(f),R(f)$, the earlier method initialized
$V$ at $f$ to the trailing secant $v=(P(f)-P(f-\Delta))/\Delta$.
Its stored values were
\[
 (v, P(f)-fv, Q(f)+R(f)/(2v_{\rm safe}), e^{2P(f)}R(f)).
\]
At a known depth the original code's clipping convention makes
those values algebraically equivalent to the observations used
here. Thus the secant restart copies $P,Q,R$ at the endpoint.
It changes the missing $V$ component. The change in backward
field history measures the performance of the restart rule.

A finite-window inverse estimates $V$ by integrating the reduced
equations back to $f-\Delta$ from a trial endpoint velocity $V_n$.
If $P_{\mathrm{trial}}(f-\Delta;V_n)$ is the resulting field, the
iteration is
\[
 V_{n+1}=V_n+
 \frac{P_{\mathrm{trial}}(f-\Delta;V_n)-P(f-\Delta)}{\Delta}.
\]
Only the specified observations are used in this iteration.
The endpoint truth velocity and earlier comparison times are
reserved for evaluation. To see why the iteration can work, write
$P_{\mathrm{trial}}(f-\Delta;V)=P(f)-\Delta V+J(V)$, where
$J(V)=\int_0^\Delta(\Delta-u)P_{\tau\tau}(f-u;V)\,du$.
For a fixed finite-dimensional discretization, a state-flow
Lipschitz bound $e^{Lu}$ and acceleration Lipschitz bound $L_F$
make $J/\Delta$ Lipschitz with constant at most
$\frac12\Delta L_F e^{L\Delta}$ on an admissible ball. An
invariant ball and a bound below one would give local convergence.
We do not use this sufficient bound as a certificate for the
simulations.

At a homogeneous polarized vacuum background,
a finite Fourier Galerkin system gives a further analytic check.
When the initial $Q$ and $Q_t$ perturbations vanish, the
linearized equation for the unmeasured $P_t$ component
is the polarized equation. If all the retained two-time
multipliers are nonzero, the derivative of the map from that
velocity to the second field snapshot is invertible.
The finite-dimensional inverse function theorem then gives
a local inverse of this reduced observation map.
The neighborhood and constants can depend on the bandwidth.

The exact true-spike family from Lim provides a local
known-answer control \cite{lim,lim2}. The corrected form for
$Q$ already occurs in Appendix A of \cite{lim2}. For fixed
time, $e^{-P(x)}$ is quadratic about the spike center.
The center can be estimated from the vertex using sampled
$P$ alone. Evaluating the two fitted quadratics there gives
the finite-window center slope and fixes the family velocity
parameter. Center values of $P$ and $Q$ then determine the
remaining family parameters. This procedure uses information
available in the original tuple and the formulas of a known
family. An off-cell center can be recovered without providing
its grid shift. The calculation is local in $x$ and does not
identify a general periodic feature as a Lim solution.

The original production example had $P=R=0$,
$V=2+1.5\cos\theta$, and $Q=\cos\theta$ at $\tau=0$.
The ordinary secant and the inverse receive identical observations
at $f-2,f$ with $f=12$ or $18$. We retain the original mask at
depth 18, while also scoring the full periodic domain.
For two states, let $d_H$ be hyperbolic distance in the $(Q,e^{-P})$
upper half-plane, $s=(V^2+e^{2P}R^2)^{1/2}$, and
$E=s^2+e^{-2\tau}(P_\theta^2+e^{2P}Q_\theta^2)$.
The score is the sampled supremum of
$\bigl(d_H^2+(\Delta s)^2+[(\Delta E)/(1+|E_{\mathrm{ref}}|)]^2\bigr)^{1/2}$.
It involves the reduced fields and their derivatives.
The comparison omits the value of the metric potential.

At $f=18$ the following scores are measured at the held-out
slice $\tau=2$. The legacy forward stencil is second order.
The two inverse columns use a fourth-order restart. The
privileged case supplies the endpoint velocity from the reference.
\begin{center}
\begin{tabular}{rccccc}
\toprule
 & \multicolumn{2}{c}{Secant, legacy stencil} &
 \multicolumn{2}{c}{Window inverse, fourth order} & Privileged \\
$N$ & Mask & Full circle & Mask & Full circle & Full circle\\
\midrule
1024 & $0.82012$ & $2.73288$ & $0.04666$ & $0.04666$ & $0.04734$\\
2048 & $0.81689$ & $3.09459$ & $0.01186$ & $0.01186$ & $0.01203$\\
4096 & $0.82215$ & $3.03402$ & $0.00298$ & $0.00298$ & $0.00302$\\
\bottomrule
\end{tabular}
\end{center}
At depth 18 the window fit residual divided by the observation
gap is below $10^{-10}$ at every tested grid. On the finest
grid the inverse gives $1.15\times10^{-3}$ and $1.55\times10^{-4}$
at $\tau=4$ and $8$. Depth 12 has no flagged component.
There the fine-grid full-circle score at $\tau=2$ falls from
$0.43207$ to $9.95\times10^{-5}$.
A same-stencil inverse is also recorded and tests inversion of its
own discretization. An analytic periodic geodesic wave-map control
provides snapshots independently of the solver. Its fourth-order
window inverse reduces the $\tau=2$, $N=1024$ score from
$0.02577$ to $1.62\times10^{-9}$.
For the two noiseless local Lim controls, the estimated half-cell
center errors were at most $1.45\times10^{-15}$ in coordinates.
Those controls test recovery within a specified family.

For a separate nonlinear check, evolve from $P=0$,
$V=2+1.4\cos\theta$, $Q=0.9\cos\theta+0.1\sin2\theta$,
and $R=0.025\cos3\theta$, with a periodic primitive of the
momentum source. Observe at $\tau=10,12$ using a fourth-order
spatial stencil. Reconstruct using the historical second-order
stencil and compare at $\tau=2$.
Apply the same reduced full-circle score.

\begin{center}
\begin{tabular}{rccc}
\toprule
$N$ & Secant & Window inverse & Exact endpoint velocity\\
\midrule
512  & $0.23690$ & $0.009585$ & $0.009876$\\
1024 & $0.23614$ & $0.002405$ & $0.002479$\\
\bottomrule
\end{tabular}
\end{center}
The matched-window score is close to that of a restart supplied
with the exact numerical endpoint velocity. On the finer grid,
adding an absolute error of $10^{-7}$ to the earlier $P$
observation in a single cosine mode gave gains of 5.17 and 8.36
in the comparison score at modes 1 and 128, respectively.
These are two directional measurements for this IVP. They
provide no infinite-band bound. The thin-feature classification
and the asymptotic regime of nonlinear unpolarized Gowdy
remain questions distinct from finite-window reconstruction.


\section{Conclusions}
In the polarized problem, the same pair of field snapshots has
different regularity costs for different targets. A rational gap can
make the field reconstruction stable at specified slices while the
missing velocity remains sensitive to high-frequency errors.
Assigning a sampling time to each band removes that additional
loss in the energy-normalized velocity.
The nonlinear comparison shows that the large error in the
original finite-window test can be removed with the same inputs
by fitting the PDE evolution across that window.
A uniform nonlinear inverse estimate near true spikes remains
open. The analytic statements concern a fixed periodic marking and
the specified finite-time observation maps.

\appendix
\section{Derivation of the scalar and metric signs}
In areal time write the base metric as $h=e^{2\sigma}(-dt^2+d\theta^2)$,
where $\sigma=-(\log t+\lambda)/4$.
The two one-dimensional warped directions have logarithmic
factors $A=(\log t+P)/2$ and $B=(\log t-P)/2$.
The base components of the four-dimensional Ricci tensor are
$R_{ab}(h)-\nabla^h_a\nabla^h_b(A+B)-A_{,a}A_{,b}-B_{,a}B_{,b}$.
Its mixed component and the sum of the two diagonal components are
\begin{align*}
 R_{t\theta}&=-\lambda_\theta/(4t)-P_tP_\theta/2,\\
 R_{tt}+R_{\theta\theta}
   &=-\lambda_t/(2t)-(P_t^2+P_\theta^2)/2 .
\end{align*}
Apply $R_{ab}=\partial_a\varphi\,\partial_b\varphi$ to obtain
\eqref{constraint}. The remaining base difference follows by
differentiating those quadratures and using the wave equations.
The two orbit equations reduce to \eqref{wave} for $P$.
The scalar d'Alembertian gives the same equation for $\varphi$.
As a check on homogeneous slices, their intrinsic scalar curvature
vanishes, so $G(n,n)=H_\theta H_x+H_xH_y+H_yH_\theta$.
With $P_\tau=p$, $\phi_\tau=v$ and $\lambda_\tau=p^2+v^2$
the expression reduces to $v^2/(4N^2)$. The vacuum example $p=1,v=0$
and the stiff isotropic example $p=0,v=\sqrt{3}$ pass this constraint.


\clearpage
\begin{thebibliography}{99}
\bibitem{aldroubi} A. Aldroubi, C. Cabrelli, U. Molter and S. Tang,
Dynamical sampling, \href{https://arxiv.org/abs/1409.8333}{arXiv:1409.8333}.
\bibitem{amorim} P. Amorim, C. Bernardi and P. G. LeFloch,
Computing Gowdy spacetimes via spectral evolution in future and past directions,
\emph{Class. Quantum Grav.} 26 (2009), 025007.
\href{https://doi.org/10.1088/0264-9381/26/2/025007}{DOI}.

\bibitem{epd} P. Amorim and P. G. LeFloch, Sharp estimates for periodic
solutions to the Euler--Poisson--Darboux equation,
\emph{Port. Math.} 65 (2008), 387--429.
\href{https://ems.press/content/serial-article-files/44605}{Full text}.

\bibitem{beyer} F. Beyer and P. G. LeFloch, Second-order hyperbolic
Fuchsian systems. Gowdy spacetimes and the Fuchsian numerical algorithm,
\href{https://arxiv.org/abs/1006.2525}{arXiv:1006.2525}.

\bibitem{cim} P. T. Chru\'sciel, J. Isenberg and V. Moncrief,
Strong cosmic censorship in polarized Gowdy spacetimes,
\emph{Class. Quantum Grav.} 7 (1990), 1671--1680.
\href{https://homepage.univie.ac.at/piotr.chrusciel/papers/chrusciel/PolarisedGowdyCQG.pdf}{Full text}.

\bibitem{snapshot} J. Christensen, F. Gonzalez, T. Kakehi and J. Wang,
The snapshot problem for the wave equation,
\emph{Advances in Mathematics} 444 (2024), 109629.
\href{https://arxiv.org/abs/2308.12208}{Preprint}.

\bibitem{li} W. Li, Scattering towards the singularity for the wave
equation and the linearized Einstein--scalar field system in Kasner spacetimes,
\href{https://arxiv.org/abs/2401.08437}{arXiv:2401.08437}.

\bibitem{lim} W. C. Lim, New explicit spike solutions---non-local
component of the generalized Mixmaster attractor,
\emph{Class. Quantum Grav.} 25 (2008), 045014.
\href{https://doi.org/10.1088/0264-9381/25/4/045014}{DOI}.

\bibitem{lim2} W. C. Lim, L. Andersson, D. Garfinkle and F. Pretorius,
Spikes in the mixmaster regime of $G_2$ cosmologies,
\emph{Phys. Rev. D} 79 (2009), 123526.
\href{https://arxiv.org/abs/0904.1546}{Preprint}.

\bibitem{rw} A. D. Rendall and M. Weaver,
Manufacture of Gowdy spacetimes with spikes,
\emph{Class. Quantum Grav.} 18 (2001), 2959--2975.
\href{https://arxiv.org/abs/gr-qc/0103102}{Preprint}.

\bibitem{szi} A. Sz\'ij\'art\'o and J. Heged\H{u}s,
Observation problems posed for the Klein--Gordon equation,
\emph{Electron. J. Qual. Theory Differ. Equ.} 2012, No. 7, 1--13.
\href{https://doi.org/10.14232/ejqtde.2012.1.7}{DOI}.

\bibitem{vest} A. Vest, Observation of vibrating systems at different time
instants, \emph{Electron. J. Qual. Theory Differ. Equ.} 2014, No. 14, 1--15.
\href{https://www.math.u-szeged.hu/ejqtde/p2411.pdf}{Full text}.

\bibitem{dlmf} NIST Digital Library of Mathematical Functions,
Chapter 10, Bessel Functions. \href{https://dlmf.nist.gov/10}{Chapter}.

\bibitem{earlier} R. Bibb, Boundary-State Sufficiency in Gowdy Cosmology:
Smooth Closure Does Not Extend Through Spikes, v1.0.
\href{https://doi.org/10.5281/zenodo.21959052}{Zenodo (2026)}.
\end{thebibliography}
\end{document}
